\section{A second Lagrangian system}
\label{secsys:section}
\label{sec:second-system}

Retain the local $e$-locus $D$ and its cohomology marking from
Proposition~\ref{def:fixed-base}, and write $f_0:X_0\to B$ for the
original projective fibration.  The geometric target is an actual nearby
manifold with two integral nef isotropic classes whose entire positive
interval lies in its K\"ahler cone.  This gives both fibrations on one
projective manifold, with the first still mapping to the original $B$.

\begin{proposition}[Two systems on one nearby model]\label{secsys:second-system}
Arbitrarily close to the original point of $D$, there is a member
$X'=X_t$ and an integral $(1,1)$-class $u$ such that $e$ and $u$ are
nef and
\[
                  q(e)=q(u)=0,\qquad q(e,u)>0,
\]
and the whole positive interval satisfies
\begin{equation}\label{secsys:geometric-target}
             \{a e+b u:a,b>0\}\subset\mathcal K_{X'}.
\end{equation}
In particular $X'$ is projective.  Its fixed-base fibration
$f'=f_t:X'\to B$ admits a second projective Lagrangian fibration
$g:X'\to C$, with $C$ normal and projective.  For every fiber $A$
over the smooth locus of $g$, the restriction $f'|_A:A\to B$ is
finite and surjective.  Thus $B$ has a finite cover by an abelian variety.
\end{proposition}

The rank-four refinement proved in
Proposition~\ref{secsys:rank-four-variation} chooses $g$ with nonzero
polarized variation when the original fibration has that property.
We first construct the two classes in
\eqref{secsys:geometric-target}.

We distinguish integral isometries from parallel-transport operators
throughout the proof.  This distinction matters in applying Torelli.
Write $\Lambda=H^2(X_0,\Z)$, equipped with $q$, and fix the connected
marked deformation component containing $X_0$.  Its orientation of maximal
positive subspaces is the one given by
$(\operatorname{Re}\sigma_0,\operatorname{Im}\sigma_0,\kappa)$, with
$\kappa$ K\"ahler.  Markings below are anchored by parallel transport
in this component.

\subsection{Realized positive three-planes}

Call an oriented positive three-plane $T\subset\Lambda_{\R}$
\emph{realized} if it is the span of the period plane and a K\"ahler
class of a marked manifold in the chosen component.  The Calabi--Yau
theorem and twistor rotation imply that every compatible oriented
two-plane $P\subset T$, together with its positively oriented orthogonal
ray in $T$, is realized once $T$ is realized.  We use period surjectivity
from the chosen marked component for the integral Beauville lattice
\cite[\S3.1 and Theorem~5.5]{HuybrechtsTorelli2012}, together with the
marked period identification
\cite[Corollaries~4.10 and~5.9]{HuybrechtsTorelli2012}.
The K\"ahler form of parallel-transport Torelli says that a Hodge
parallel-transport isometry carrying a K\"ahler class to a K\"ahler
class is induced by an isomorphism
\cite[Theorem~1.3(2) and \S3.2]{Markman2011}.  The global Torelli input
in that proof is the preceding marked period identification.
These are the marked-component statements used below.
The twistor description is also recalled in
\cite[Section~5.1]{AmerikVerbitsky2015}.

\begin{lemma}\label{secsys:metric-periods}
Realized positive three-planes form an open subset.  In rank greater
than four, a three-plane fails to be realized precisely when it is
orthogonal to a rational negative MBM class of the marked component.
In rank four, the weaker assertion that every nonrealized three-plane
has a rational negative orthogonal line suffices and holds.
At a fixed oriented period, K\"ahler cones of marked models in the
component are either equal or disjoint.
\end{lemma}

\begin{proof}
Given a realized pair $(P,\kappa)$, nearby planes $T'$ contain oriented
period planes $P'$ close to $P$ and orthogonal positive vectors
$\kappa'$ close to $\kappa$.  Local Torelli realizes $P'$ by a small
deformation of that manifold.  Openness of the relative K\"ahler cone
makes $\kappa'$ K\"ahler, proving openness.

Recall the two parts of the MBM theorem used here: the negative rational
wall classes are invariant under deformation that keeps them of type
$(1,1)$, and the K\"ahler cone is a component of the positive cone
after removing their orthogonal hyperplanes
\cite[Theorems~1.17 and~1.19]{AmerikVerbitsky2015}.
No bound on their squares is used.  A realized plane cannot be
orthogonal to such a class: that class would be of type $(1,1)$ and
orthogonal to a K\"ahler class.  Conversely, suppose $T$ has no such
orthogonality.  Choose a generic oriented plane $P\subset T$, so that
every rational class orthogonal to $P$ is already orthogonal to $T$.
Indeed, each rational vector not orthogonal to $T$ excludes a proper
closed subset of the sphere of planes in $T$, and there are only
countably many such vectors.  Period surjectivity gives a marked
manifold with period $P$.  Its rational $(1,1)$ classes include no
MBM class.  The cone theorem therefore identifies its K\"ahler cone
with the whole positive component.  The orthogonal ray to $P$ in $T$,
with sign fixed by the orientation, is in that component.  This realizes
$T$.

For rank four one can avoid even the MBM formulation in this step.
If the negative line $T^\perp$ is irrational, the same generic choice
of $P$ has zero Picard lattice.  Its K\"ahler cone is the whole positive
component, so $T$ is realized.  Thus a nonrealized plane has rational
negative normal line.

Finally, an intersection between the K\"ahler cones of two marked models
with the same period gives a Hodge parallel-transport isometry carrying
a K\"ahler class to a K\"ahler class.  Torelli makes it the isometry of
an isomorphism.  The cones then agree.
\end{proof}

Choose an integral ample class $h$ on $X_0$.  A first arbitrarily small
deformation preserving both $h$ and $e$ makes the real period plane
$P_*$ rational: rational positive two-planes are dense in the rational
space $\langle e,h\rangle^\perp$.  In rank four this plane is already
rational, and we take the original point.  Ampleness of $h$ persists, and
Proposition~\ref{def:fixed-base} preserves the fibration of class $e$.
Denote this projective member and its fibration by $X_*,f_*$.
Set
\begin{equation}\label{secsys:four-space}
 U=P_*\oplus\langle e,h\rangle_{\Q},
 \qquad \operatorname{sign}(q|_U)=(3,1).
\end{equation}
Here and below a rational plane is also used for its real span when
discussing signs and periods.  The space
$P_*+\R h$ is realized.  Consequently $U^\perp$ contains no MBM
class that could obstruct realization.

Let $\mathbb H(U)$ be the hyperbolic three-space of negative lines in
$U_{\R}$.  A point $[v]$ corresponds to the positive three-plane
$v^\perp\cap U_{\R}$, with its fixed orientation.  Let
$J\subset\mathbb H(U)$ be the set of nonrealized planes.  It is closed
by Lemma~\ref{secsys:metric-periods}.  Moreover,
\begin{equation}\label{secsys:rational-bad-set}
             J\subset\mathbb H(U)(\Q).
\end{equation}
For rank greater than four, project an obstructing rational MBM class
orthogonally to $U$.  The projection is rational and nonzero, since
$U^\perp$ contains no such class; it is the negative normal line in
$U$.  In rank four use the last assertion about nonrealized planes in
Lemma~\ref{secsys:metric-periods}.  Any parallel-transport isometry
preserving $U$ preserves $J$.

Choose the component of the negative cone whose null boundary contains
the ray $\R_{>0}e$.  We call its nonzero vectors, including boundary
vectors, future vectors.  Distinct future null vectors have negative
pairing.  The other null line in $\langle e,h\rangle$ has a rational
future representative
\begin{equation}\label{secsys:other-cusp}
 w_*=e-\frac{2q(e,h)}{q(h)}h.
\end{equation}
The negative lines in a Lorentz plane with future null boundary vectors
$v,w$ form the geodesic
$\{[s v+t w]:s,t>0\}$; we denote it by $(v,w)$.

\subsection{Null rotations from actual algebraic families}

\begin{lemma}\label{secsys:actual-monodromy}
Let $a\in U_{\Q}$ be nonzero and null.  Every rational null rotation
fixing $a$, extended as the identity on $U^\perp$, has a positive power
that is an integral parallel-transport operator in the chosen marked
component.
\end{lemma}

\begin{proof}
Choose $b\in U_{\Q}$ with $q(a,b)=1$ and $q(b)=0$; subtracting
$q(b_0)a/2$ from any $b_0$ paired to $1$ with $a$ gives such a vector.
Let $V=\langle a,b\rangle^\perp\cap U$, a positive rational plane.
For $x\in V_{\Q}$ the null rotation is
\begin{equation}\label{secsys:null-rotation}
 \begin{split}
 T_{a,x}a&=a,\\
 T_{a,x}b&=b+x-\tfrac12 q(x)a,\\
 T_{a,x}v&=v-q(v,x)a\quad(v\in V),
 \end{split}
 \qquad T_{a,x}|_{U^\perp}=\id.
\end{equation}
Direct pairing verifies that it preserves $q$, and substitution gives
$T_{a,x}T_{a,y}=T_{a,x+y}$.  Thus $T_{a,x}^m=T_{a,mx}$.
Its matrix entries in an integral basis are polynomials in $m$ of degree
at most two with rational coefficients and constant term the identity.
For sufficiently divisible $m$, both $T_{a,mx}$ and its inverse are
integral.  This accounts also for the finite gluing between the
sublattices in $U$ and $U^\perp$.

For $x\ne0$, choose a rational positive vector
$\ell\in V\cap x^\perp$.  The rotation fixes $\ell$.  Replace it by
a primitive integral vector on the same line.  A generic period in
$\ell^\perp$ has Picard lattice $\Z\ell$.  Period surjectivity and
the cone theorem give an ample marked deformation with polarization
$\ell$ or $-\ell$; choose the ample sign.  This sign does not alter
the stabilizer of $\ell$.

We now use the prescribed-polarization conclusion of
\cite[Lemma~4.2]{StrongSYZ2026}.  It applies to this primitive positive
class $\ell$, ample on a manifold in the chosen marked component.
Write $D_\ell$ for the component of its type-IV period domain containing
that ample period, and put
\[
 G_\ell(\Z)=\{g\in O(\Lambda,q):g\ell=\ell,
                                      \ gD_\ell=D_\ell\}.
\]
The cited lemma supplies a smooth projective family of IHS manifolds
over a smooth connected quasi-projective variety $S$, with a relative
polarization of class $k\ell$ for some $k>0$ and a reference marking
obtained by parallel transport from $X_0$.  It also supplies a
torsion-free finite-index subgroup $\Gamma\subset G_\ell(\Z)$
containing the family monodromy and an algebraic period map
\[
           \Phi:S\longrightarrow Q_\ell:=\Gamma\backslash D_\ell.
\]
Here $Q_\ell$ is smooth and quasi-projective, the map
$D_\ell\to Q_\ell$ is a covering with free $\Gamma$-action, and a
local marked lift of $\Phi$ is submersive.  The type-IV domain
$D_\ell$ is contractible, including in rank four, so this covering
identifies $\pi_1(Q_\ell)$ with $\Gamma$.

The remaining finite-index assertion concerns the image of this
actual family monodromy, and we prove it here.  The submersive local
lift makes the algebraic map $\Phi$ dominant.

Choose an irreducible algebraic slice $S'\subset S$ of dimension
$\dim Q_\ell$ dominating $Q_\ell$.  For example, intersect with
general hyperplanes through a smooth submersive point, transverse to
its period fiber.  After normalization and restriction, this slice is
smooth and generically finite over $Q_\ell$.  There is a nonempty
Zariski-open $Q^\circ\subset Q_\ell$ over which its restriction is
a nonempty finite \'{e}tale cover: generic finiteness gives a finite
open after removing the boundary in a proper compactification, and
characteristic zero allows removal of the branch locus.
Therefore its fundamental-group image has finite index in
$\pi_1(Q^\circ)$.  The map
$\pi_1(Q^\circ)\to\pi_1(Q_\ell)$ is surjective, since a loop in a
complex manifold can be moved off a proper algebraic subset.
It follows that
\begin{equation}\label{secsys:finite-index-image}
       [\Gamma:\Phi_*\pi_1(S)]<\infty.
\end{equation}
The period map on the universal cover of $S$, based at this reference
marking, is equivariant for the actual second-cohomology monodromy.
Because $D_\ell\to Q_\ell$ is a
covering with free action, the deck transformation induced by each
loop is exactly that cohomology operator.  Hence the subgroup in
\eqref{secsys:finite-index-image} consists of actual parallel transports,
with the original marking retained.

Our rotation fixes $\ell$; after taking an integral power, a further
power preserves $D_\ell$ and therefore lies in $G_\ell(\Z)$.
Since $\Gamma$ has finite index in $G_\ell(\Z)$, a further power
lies in $\Gamma$; by \eqref{secsys:finite-index-image}, one more
positive power lies in the family monodromy.  This uses only that every element has a
positive power in any finite-index subgroup; normality of the subgroup
is unnecessary.  Transport to the original reference fiber finishes
the proof.  The zero parameter gives the identity.
\end{proof}

\subsection{An integral loxodromic operator with irrational axis}

\begin{lemma}\label{secsys:loxodromic}
There is a parallel-transport operator $\gamma$ preserving $U$ and
acting identically on $U^\perp$ such that its characteristic polynomial
on $U$ is irreducible over $\Q$.  Its real eigenvalues are
$\rho,\rho^{-1}$ with $\rho>1$, and its other two eigenvalues are a
nonreal conjugate pair on the unit circle.  Its attracting and repelling
future null lines $a_+,a_-$ can both be placed in any prescribed
neighborhood of $[w_*]$ in the null boundary.
\end{lemma}

\begin{proof}
Choose a rational null partner $b$ of $e$ on the line of $w_*$ with
$q(e,b)=1$, and a rational orthogonal basis $v_1,v_2$ of their positive
perpendicular plane.  Write $q(v_1)=A>0$, $q(v_2)=B>0$.
For a rational number $t\ne0$, set $x=v_1$, $y=v_1+t v_2$.
Choose $t$ so that
\begin{equation}\label{secsys:phase-choice}
                  \frac{2(A-Bt^2)}{A+Bt^2}\notin\Z.
\end{equation}
This is possible: the expression lies in $(-2,2)$ and assumes each
of the three possible integral values at only finitely many real $t$.
By Lemma~\ref{secsys:actual-monodromy}, for some positive integers
$m_1,m_2$ the opposite rotations $T_{e,m_1x}$ and $T_{b,m_2y}$ are
integral monodromies.

For completeness, the spectral calculation can be performed in a real
orthonormal transverse basis without any rationality assumption on
$\sqrt A$ or $\sqrt B$.  Identify
$a e+b' b+c v_1+d v_2$ with the Hermitian matrix
\[
 \begin{pmatrix}
 \sqrt2 a&\sqrt A c+i\sqrt B d\\
 \sqrt A c-i\sqrt B d&-\sqrt2 b'
 \end{pmatrix}.
\]
Its determinant is the negative of $q$.  The action
$H\mapsto MHM^*$ of $SL_2(\C)$ induces $SO^+(3,1)$.
Formula~\eqref{secsys:null-rotation} corresponds, for the first
rotation, to an upper triangular unipotent with parameter
$-m_1\sqrt A/\sqrt2$; the opposite rotation corresponds to a lower
triangular unipotent with parameter
$m_2(\sqrt A-i t\sqrt B)/\sqrt2$.
Put
\[
            \gamma_k=T_{e,km_1x}T_{b,km_2y}.
\]
Its spin trace is
\begin{equation}\label{secsys:spin-trace}
 \tau_k=2+k^2m_1m_2\zeta,
 \qquad \zeta=-\tfrac12(A-i t\sqrt{AB}).
\end{equation}
Let $\lambda_k,\lambda_k^{-1}$ be its spin eigenvalues, choosing
$|\lambda_k|>1$.  For large $k$,
$\lambda_k/\tau_k\to1$.  The four eigenvalues on $U_{\R}$ are
\[
 |\lambda_k|^2,\quad |\lambda_k|^{-2},\quad
 \lambda_k/\overline{\lambda_k},\quad
 \overline{\lambda_k}/\lambda_k.
\]
The sum of the last pair converges to
$2\cos(2\arg\zeta)=2(A-Bt^2)/(A+Bt^2)$.
By \eqref{secsys:phase-choice}, it is nonintegral for all sufficiently
large $k$.  In particular, this conjugate pair is nonreal and neither
eigenvalue is $1$ or $-1$.

The characteristic polynomial on $U$ is monic with rational
coefficients, and its roots are algebraic integers because
$\gamma_k$ preserves $\Lambda$.  It therefore belongs to $\Z[z]$.
It is reciprocal, has constant term $1$, and has no rational root:
a rational root would be $1$ or $-1$.  If it factored into two
quadratics over $\Q$, real coefficients would force the nonreal pair
to belong to one quadratic.  That monic factor has integral
coefficients, so the sum of the pair would be an integer, a
contradiction.  Thus it is irreducible.  Fix a sufficiently large $k$.

Finally, a nontrivial null rotation fixing $[w_*]$ has powers converging
to $[w_*]$ on every other point of the null boundary.  In boundary
coordinates with $[w_*]$ at infinity these powers are translations
$z\mapsto z+jc$, with $c\ne0$.  Such a rotation has an integral
monodromy power by Lemma~\ref{secsys:actual-monodromy}.  Neither
eigenline of $\gamma_k$ is rational, by irreducibility, so neither is
$[w_*]$.  Conjugating $\gamma_k$ by a sufficiently high power of this
rotation brings both eigenlines as close to $[w_*]$ as desired and
preserves all the asserted properties.  Call the conjugate $\gamma$.
\end{proof}

The axis $(a_-,a_+)$ contains no rational negative line.  Indeed, a
nonzero rational vector in its invariant real Lorentz plane would
have its entire $\gamma$-orbit in that plane.  The rational span of
the orbit would be a nonzero proper rational invariant subspace of
$U$, contradicting irreducibility.  The two geodesics $(e,a_\pm)$
also contain no rational negative line.  Otherwise such a line,
together with rational $e$, would define a rational Lorentz plane.
A rational binary quadratic form with one rational null line has
both null lines rational: if $q(e,v)\ne0$, its second null vector is
$v-q(v)e/(2q(e,v))$.  The second line would be $a_\pm$, contrary
to irreducibility.  Equation~\eqref{secsys:rational-bad-set} gives
\begin{equation}\label{secsys:good-limit-geodesics}
 J\cap(a_-,a_+)=J\cap(e,a_+)=J\cap(e,a_-)=\varnothing.
\end{equation}

\subsection{Uniform tails and the entire K\"ahler interval}

For a future null line $[w]$ near $[w_*]$, choose its representative
with $q(e,w)=-1$ and put
\[
              P(e,w)=\langle e,w\rangle^\perp\cap U_{\R},
\]
oriented near $P_*$.  This is a nearby $e$-period.  It gives an actual
small deformation $X_w$ of $X_*$ within $D$, with the fibration $f_w:X_w\to B$ by
Proposition~\ref{def:fixed-base}.  The orthogonal projection of $h$ to
$\langle e,w\rangle$ has the form
\begin{equation}\label{secsys:projected-kahler}
 h_w=\alpha(w)e-\beta w,\qquad
 \alpha(w)=-q(h,w)>0,\quad \beta=q(h,e)>0.
\end{equation}
It tends to $h$ as $[w]\to[w_*]$, so it is K\"ahler on $X_w$
after restricting to a small neighborhood.  As $e$ is nef on $X_w$,
every $h_w+t e$, $t\ge0$, is K\"ahler.  Its orthogonal negative
line in $\langle e,w\rangle$ is
$[(\alpha(w)+t)e+\beta w]$.  On a smaller neighborhood $\mathcal V$
of $[w_*]$, the continuous function $\alpha(w)/\beta$ is bounded
above by some $R$.  Thus
\begin{equation}\label{secsys:uniform-tail}
 [s e+w]\notin J\qquad
 ([w]\in\mathcal V,\ q(e,w)=-1,\ s\ge R).
\end{equation}
This is the required uniformity near the rational endpoint $e$.
Choose the conjugation in Lemma~\ref{secsys:loxodromic} so that both
$a_+$ and $a_-$ belong to $\mathcal V$.

\begin{lemma}\label{secsys:whole-geodesic}
For all sufficiently large integers $N$,
\[
                         (e,\gamma^N e)\cap J=\varnothing.
\]
\end{lemma}

\begin{proof}
Choose future eigenvectors for $a_+,a_-$ with eigenvalues
$\rho,\rho^{-1}$, respectively.  The remaining real invariant
plane is positive definite and $\gamma$ acts on it by a rotation.
Since $e$ is a future null vector different from both eigenlines,
its coefficients along both future eigenvectors are positive.
Indeed, pairing its eigenspace decomposition with either eigenvector
shows the coefficient of the other to be positive.  It follows that
\begin{equation}\label{secsys:endpoint-limits}
 \rho^{-N}\gamma^Ne\longrightarrow v_+,
 \qquad
 \rho^{-N}\gamma^{-N}e\longrightarrow v_-,
\end{equation}
where $v_+,v_-$ are positive representatives of $a_+,a_-$.

If the assertion failed, choose $N\to\infty$ and
$z_N=[s_Ne+\gamma^Ne]\in J$, with $s_N>0$.  We first claim
\begin{equation}\label{secsys:balance-bounds}
             s_N\rho^{-N}\longrightarrow0,
             \qquad s_N\rho^N\longrightarrow\infty.
\end{equation}
To prove the first limit, normalize $\gamma^Ne$ by
$c_N=-q(e,\gamma^Ne)$; by \eqref{secsys:endpoint-limits},
$c_N/\rho^N$ tends to a positive constant, and
$[\gamma^Ne]\to a_+\in\mathcal V$.
The tail exclusion~\eqref{secsys:uniform-tail} bounds $s_N/c_N$
above.  If a subsequence of $s_N/\rho^N$ had a positive limit,
$z_N$ would converge to an interior point of $(e,a_+)$.
Closedness of $J$ would put that point in $J$, contradicting
\eqref{secsys:good-limit-geodesics}.  This proves the first limit.
Apply the same argument to
\[
 \gamma^{-N}z_N
       =[s_N\gamma^{-N}e+e]
       =[s_N^{-1}e+\gamma^{-N}e]\in J.
\]
The uniform tail at $e$ now uses $a_-\in\mathcal V$ and gives
$s_N^{-1}\rho^{-N}\to0$, which is the second limit.

Choose integers $k_N$ such that
\[
               1\le s_N\rho^{2k_N-N}<\rho^2.
\]
The two limits in \eqref{secsys:balance-bounds} imply
$k_N\to\infty$ and $N-k_N\to\infty$.
Passing to a subsequence, suppose
$s_N\rho^{2k_N-N}\to r\in[1,\rho^2]$.
Apply $\gamma^{-k_N}$ to $z_N$ and divide its representative by
$\rho^{N-k_N}$.  Equation~\eqref{secsys:endpoint-limits} gives
\[
 \frac{s_N\gamma^{-k_N}e+\gamma^{N-k_N}e}{\rho^{N-k_N}}
   =s_N\rho^{2k_N-N}
       \frac{\gamma^{-k_N}e}{\rho^{k_N}}
       +\frac{\gamma^{N-k_N}e}{\rho^{N-k_N}}
   \longrightarrow r v_-+v_+.
\]
The limit has negative square, so its line is an interior point of
the axis $(a_-,a_+)$.  Invariance and closedness of $J$ put that
line in $J$, contradicting \eqref{secsys:good-limit-geodesics}.
\end{proof}

Take such an $N$ and let $X_N$ be the actual nearby $e$-deformation
with first fibration $f_N:X_N\to B$ and period
$P_N=P(e,\gamma^Ne)$, oriented near $P_*$.
The choice of $a_+$ close to $[w_*]$, followed by large $N$, makes
this deformation as small as prescribed.  Put
\begin{equation}\label{secsys:second-class}
 u_N=-\gamma^Ne,\qquad
 \mathcal C_N=\{a e+b u_N:a,b>0\}.
\end{equation}
The minus sign is essential: $e$ and $\gamma^Ne$ are future null
vectors, so $q(e,\gamma^Ne)<0$ and $q(e,u_N)>0$.
The cone $\mathcal C_N$ is a positive component in the Lorentz plane
$\langle e,u_N\rangle$.  Its positive rays correspond, by taking
orthogonal negative lines, to the geodesic $(e,\gamma^Ne)$.
Lemma~\ref{secsys:whole-geodesic} and twistor rotation show that every
ray of $\mathcal C_N$ is K\"ahler on some marked model of period $P_N$.

All those rays are K\"ahler on $X_N$ itself.  Indeed, their projectivized
interval is connected.  Intersections with the K\"ahler cones of marked
models give an open cover by pairwise disjoint sets, by
Lemma~\ref{secsys:metric-periods}.  Such a cover has just one nonempty
member.  The class $h_w$ of \eqref{secsys:projected-kahler}, for
$[w]=[\gamma^Ne]$, belongs to this interval and is K\"ahler on $X_N$.
Thus
\begin{equation}\label{secsys:actual-cone}
                         \mathcal C_N\subset\mathcal K_{X_N}.
\end{equation}
In particular $e+u_N$ is an integral K\"ahler class, making $X_N$
projective, and $u_N$ is a nonzero integral nef isotropic $(1,1)$
class on this very manifold.  It is the Chern class of a line bundle:
the exponential sequence and $H^1(X_N,\cO_{X_N})=0$ identify
$\Pic(X_N)$ with the integral $(1,1)$ classes.  Applying
Theorem~\ref{thm:strong-syz}, including Stein factorization, gives
the second Lagrangian fibration $g:X_N\to C$ with normal projective
base $C$.  Both $f_N$ and $g$ are projective morphisms: their graphs
are closed analytic subvarieties of projective products, hence
algebraic, and the corresponding algebraic maps are projective.

\subsection{Finite maps from its abelian fibers}

Let $A$ be a fiber over the smooth locus of $g$.  It is an abelian
variety after choosing an origin.  A positive multiple of $u_N$ is
$g^*c_1(H_C)$ for an ample line bundle $H_C$ on $C$, so the cohomology
class of $A$ is a positive rational multiple of $u_N^n$.  To see the
constant, cut $C$ by $n$ general hyperplanes in a sufficiently high
multiple of $H_C$; their transverse intersection is a finite collection
of regular points, and its inverse image is the corresponding collection
of fibers.  All regular fibers have the same cohomology class by
parallel transport over the connected smooth locus.

Polarizing Fujiki gives the exact positive quantity
\begin{equation}\label{secsys:mixed-fujiki}
 \binom{2n}{n}\int_{X_N}e^n u_N^n
       =c_{X_N}\,2^n q(e,u_N)^n>0.
\end{equation}
Therefore the nef line bundle $f_N^*H|_A$ has positive top
self-intersection.  Such a line bundle on an abelian variety is ample.
Indeed, average representatives of its class by translations.  Nefness
makes the resulting invariant Hermitian form positive semidefinite,
and positive top intersection makes its determinant positive.  All
its eigenvalues are then positive, so the line bundle has a positive
curvature representative and is ample.

The morphism $f_N|_A$ has no positive-dimensional fiber, since a curve
in such a fiber would have degree zero against this ample line bundle.
It is proper, hence finite.  Its image is closed of dimension
$\dim A=n=\dim B$ and therefore equals $B$.  This proves all assertions
of Proposition~\ref{secsys:second-system}.  In particular we have obtained
\begin{equation}\label{secsys:abelian-cover}
                         p:A\longrightarrow B
\end{equation}
as a finite surjective morphism from an abelian variety.

\subsection{Nonzero variation in rank four}

\begin{proposition}[The rank-four variation refinement]
\label{secsys:rank-four-variation}
If $b_2(X_0)=4$ and the polarized family of smooth fibers of $f_0$
has nonzero variation, the choices in
Proposition~\ref{secsys:second-system} can be made so that $g$ also
has nonzero polarized variation.
\end{proposition}

\begin{proof}
Nonzero variation means that, at some regular fiber, the derivative
in a base direction of its weight-one period map is nonzero.  This
condition persists under small deformations of the fibration: on a
small smooth product neighborhood in the family supplied by
Proposition~\ref{def:fixed-base}, the relative Hodge bundle and that
derivative vary holomorphically, so a chosen nonzero matrix entry
stays nonzero.

The fiber polarization can also be transported.  Its class is the
restriction of an integral ample class on $X_0$.  By
Lemma~\ref{def:restriction-kernel} and topological transport, the
restriction image on a nearby smooth fiber is a rational
one-dimensional space with kernel $e^\perp$.  The restriction of a
K\"ahler class spans its real form, since it pairs positively with $e$,
so the transported integral class
remains of type $(1,1)$.  Positivity on the smooth torus is open.
This gives a polarized family near the chosen fiber, and the same
argument applies at the conjugate fibration.

In rank four the preliminary point was $X_*=X_0$.  Complex
conjugation of its hyperk\"ahler complex structure reverses the
period-plane orientation and the K\"ahler sign.  The resulting marked
manifold $\overline X_*$ has period $\overline P_*$, ample class $-h$,
and conjugate fibration
$\overline f_*:\overline X_*\to\overline B$ with fiber class $-e$.
It lies in the same marked deformation component: a half-circle of
the twistor sphere joins the complex structure to its negative.
Nonzero variation of $\overline f_*$ is equivalent to that of $f_*$.

We choose the endpoint neighborhood in
Lemma~\ref{secsys:loxodromic} small enough for the following comparison
with this conjugate fibration.  Under the actual parallel transport
$\gamma^{-N}$, the underlying plane of $P_N$ becomes
\[
           \langle \gamma^{-N}e,e\rangle^\perp\cap U_{\R}.
\]
It approaches an unoriented plane near $P_*$, since
$[\gamma^{-N}e]\to a_-$ and $a_-$ can be chosen arbitrarily close to
$[w_*]$.  The transported positive cone is
\begin{equation}\label{secsys:opposite-cone}
        \gamma^{-N}\mathcal C_N
             =\{a\gamma^{-N}e-b e:a,b>0\}.
\end{equation}
By \eqref{secsys:other-cusp}, this cone contains a vector close to
$-h$.  Preservation of the orientation of maximal positive
three-spaces therefore makes the transported oriented period close
to $\overline P_*$.

Choose the endpoint neighborhood sufficiently small, and then $N$
sufficiently large, that this period lies in a neighborhood where the
$(-e)$-deformation of $\overline f_*$ exists and has nonzero variation.
Let $Y_N$ be that actual small deformation of $\overline X_*$ with
period $\gamma^{-N}P_N$.  It has a K\"ahler class near $-h$, in the
component \eqref{secsys:opposite-cone}.  In rank four,
$H^{1,1}(X_N,\R)=\langle e,u_N\rangle_{\R}$, so
\eqref{secsys:actual-cone} says that the entire positive component
of $X_N$ is K\"ahler.  The Hodge parallel-transport isometry between
$Y_N$ and $X_N$ induced by $\gamma^N$ consequently carries a
K\"ahler class to a K\"ahler class.  Torelli identifies these
manifolds by an isomorphism inducing that isometry.  The class $-e$
on $Y_N$ becomes $u_N$ on $X_N$.

Line bundles on an IHS manifold are determined by their Chern
classes.  Thus the transported fibration and the fibration obtained
from $u_N$ have the same connected fibers: positive generated powers
define the same Stein factorization.  It follows that $g$ has nonzero
polarized variation.  The isometry in this comparison is the actual
parallel transport constructed in Lemma~\ref{secsys:actual-monodromy}.
\end{proof}
